\subsection{Indexed Families and Products}

\begin{definition} \label{def:indexed-families-unions-intersections} \hlt{Families, Indexed Families, Unions, and Intersections}\\
Let $X$ be a set and let $I$ be an index set.
\begin{enumerate}[label=(\roman*)]
\setlength{\itemsep}{0pt}
\item \hlt{Family}: A set whose members are sets, also called a \hlt{collection} of sets.
\item \hlt{Indexed family}: An assignment $i\mapsto A_i$ over $I$, written $(A_i)_{i\in I}$ or $\{A_i\}_{i\in I}$.
It is an indexed family in $X$ iff $(\forall i\in I)\,A_i\subseteq X$.
\item \hlt{Indexed union}: For an indexed family $\{A_i\}_{i\in I}$ in $X$, $\bigcup\limits_{i\in I}A_i\defeq\{x\in X:(\exists i\in I)\,x\in A_i\}$.
\item \hlt{Indexed intersection}: For an indexed family $\{A_i\}_{i\in I}$ in $X$, $\bigcap\limits_{i\in I}A_i\defeq\{x\in X:(\forall i\in I)\,x\in A_i\}$.
\end{enumerate}
\end{definition}

\begin{theorem} \label{thm:indexed-de-morgan-laws} \hlt{Arbitrary De Morgan Laws}\\
Let $X$ and $I$ be sets, and let $\{E_i\}_{i\in I}$ be an indexed family in $X$.
Take all complements relative to $X$.
\begin{enumerate}[label=(\roman*)]
\setlength{\itemsep}{0pt}
\item $\left(\bigcup\limits_{i\in I}E_i\right)^c=\bigcap\limits_{i\in I}E_i^c$.
\item $\left(\bigcap\limits_{i\in I}E_i\right)^c=\bigcup\limits_{i\in I}E_i^c$.
\end{enumerate}
\begin{proof}
Fix an arbitrary $x\in X$.
Extensionality reduces both identities to pointwise membership equivalences.
\begin{enumerate}[label=(\roman*)]
\setlength{\itemsep}{0pt}
\item $x\in\left(\bigcup\limits_{i\in I}E_i\right)^c\iff(\forall i\in I)\,x\notin E_i\iff(\forall i\in I)\,x\in E_i^c\iff x\in\bigcap\limits_{i\in I}E_i^c$.
\item $x\in\left(\bigcap\limits_{i\in I}E_i\right)^c\iff(\exists i\in I)\,x\notin E_i\iff(\exists i\in I)\,x\in E_i^c\iff x\in\bigcup\limits_{i\in I}E_i^c$.
\end{enumerate}
\hlt{Structural Commentary}: The proof checks both identities pointwise and applies quantifier negation.\\
\hlt{Key Insight}: Complement exchanges existential and universal membership conditions.
\end{proof}
\end{theorem}

\begin{definition} \label{def:disjoint-families-and-parametrizations} \hlt{Disjoint Families and Parametrizations}\\
Let $I$ be an index set and let $\mathcal{F}$ be a family of sets.
Let $\{A_i\}_{i\in I}$ be an indexed family of sets.
\begin{enumerate}[label=(\roman*)]
\setlength{\itemsep}{0pt}
\item \hlt{Pairwise disjoint}: $(\forall E,G\in\mathcal{F})\,(E\neq G\implies E\cap G=\emptyset)$.
\item \hlt{Parametrization}: $\{A_i\}_{i\in I}$ parametrizes $\mathcal{F}$ iff $\mathcal{F}=\{A_i:i\in I\}$.
\end{enumerate}
\end{definition}

\begin{definition} \label{def:indexed-cartesian-product} \hlt{Cartesian Product of an Indexed Family}\\
Let $I$ be an index set and let $\{E_i\}_{i\in I}$ be an indexed family of sets.
\begin{enumerate}[label=(\roman*)]
\setlength{\itemsep}{0pt}
\item \hlt{Product}: $\prod\limits_{i\in I}E_i\defeq\{(x_i)_{i\in I}:(\forall i\in I)\,x_i\in E_i\}$.
\item \hlt{Binary case}: If $I=\{1,2\}$, identify $(x_i)_{i\in I}$ with the ordered pair $(x_1,x_2)$.
Under this identification, $\prod\limits_{i\in I}E_i=E_1\times E_2$.
\item \hlt{Constant family}: If $E$ is a set and $(\forall i\in I)\,E_i=E$, write $E^I\defeq\prod\limits_{i\in I}E_i$.
Thus $E^I$ is the set of $I$-indexed families in $E$.
\item \hlt{Dependence}: Different parametrizations of the same family can yield different products.
\end{enumerate}
\end{definition}
